Something had changed in Daenerys, and Jon could not name what.

It was not anything he could have pointed to, if Sansa or Davos had asked him to explain it — and he half-suspected, some mornings, that he was imagining it, inventing a distance where there was only the ordinary exhaustion of a retreat that had ground all of them down to something thinner and colder than they’d been three weeks past. But he did not think he was imagining it. He had come to know her too well across a voyage and a war, and he knew the difference between the silence of a woman worn tired and the silence of a woman watching him.

She watched him now. In council, when he spoke, he would feel her eyes on him a beat too long, weighing something in his words he could not identify. She was not cold, exactly — she still sought him out, still shared his fire on the worst nights, still let him see the small unguarded version of herself that no one else in the column was permitted. But there was a new carefulness in her, a held-back quality, as though some part of her had stepped back to a distance she had not kept from him since the earliest days aboard the ship north, before either of them had learned to trust the other.

He told himself it was Rhaegal. She had lost a child of her body to the dead, watched it grounded and killed and knew it would rise against her, and a grief like that changed a person, closed doors in them that took time to open again. That was surely all it was. He told himself so several times, in the long grey days of the march, and did not quite believe it, and could not have said why.

He did not ask her. He understood, without letting himself examine the understanding too closely, that some part of him was afraid of the answer — and a man who was afraid of an answer, Jon had learned, was often a man who already half-knew it and preferred the not-knowing.

\sceneBreak

The rearguard question came to him on the twenty-fourth day, and it was the kind of question Jon had come, these past weeks, to dread more than any battle.

The column had reached a stretch of the kingsroad where the land narrowed between a frozen river and a ridge of high broken ground — good country to be caught in, if a man had to be caught somewhere, because the narrowness meant a small force could hold the road against a larger one for a time. And a small force was going to have to, because the dead were gaining again, the way they gained every few days, tireless where the living tired, and the column’s wounded and its remaining smallfolk could not move fast enough to outrun them without someone buying the time.

“Two hundred men,” Davos said, laying it out plain, because Davos always laid things out plain. “Two hundred, holding the narrows, buys the column a full day’s march. Maybe more. Enough to reach the next defensible ground, where we can make a proper stand instead of being caught strung out along open road.”

“And the two hundred?”

Davos did not answer that directly, which was answer enough. Two hundred men holding a narrows against the dead did not come back. They held until they were overrun, and bought what time their dying could buy, and the column marched on without them.

Jon looked at the map, and understood that the decision was his, and understood, too, the particular shape of what was being asked of him — not to send men into a battle they might survive, but to choose two hundred living men and name them dead, deliberately, for the arithmetic of it. He had made hard calls before. He had never quite made this one, this cold, this clear-eyed a spending of lives that could not be spent any other way.

“Volunteers,” he said, at last. “I’ll not conscript men to a death like that. If two hundred won’t stand willingly, we find another way, whatever it costs the rest of the column.”

“And if two hundred do stand willingly?” Davos asked, quiet. “You’ll let them?”

Jon was silent a long moment. He thought of Theon, holding the godswood gap. He thought of Tormund, holding the southern gate. He thought of every man who had already bought this retreat its miles with his own life, and understood that refusing to let two hundred more do the same, out of some squeamishness about naming the cost aloud, would only mean the whole column died strung out on open road instead — that mercy here, the easy kind, the kind that could not bear to spend anyone, would kill everyone.

“I’ll let them,” he said. “And I’ll ask for them myself, to their faces, rather than send Davos or another man to do it for me. If I’m going to spend two hundred lives, I’ll not hide from the men I’m spending. They’ll hear it from me, and see that I know exactly what I’m asking, and I’ll stand there and take whatever they say back to me for the asking.”

Davos looked at him a long moment, something in his weathered face that was not quite approval and not quite grief. “That’ll cost you,” he said. “Standing in front of them. Asking it to their faces. Easier to send the order down the line and never see them again.”

“Aye,” Jon said. “It’ll cost me. It should. A man who spends lives cheaply enough that it stops costing him is a man who shouldn’t be trusted to spend them at all.”

\sceneBreak

He asked them himself, that evening, two hundred men drawn up in the cold, and told them plainly what he needed and what it would cost them, and did not dress it in any glory, because they deserved better than glory — they deserved the truth, which was that they would die, and that their dying would save the column, and that he would carry the weight of having asked it for the rest of whatever life the war left him.

More than two hundred stepped forward. That was the thing that undid him, standing there — that he’d braced to plead for volunteers and instead had to turn men away, choose among the willing, tell some brave men to march on and live because he only needed two hundred and had been offered three. He shook the hand of every man who stayed. He learned as many of their names as the cold and the hour permitted. He did not let himself look away from any of them, because he had decided that a king who spent lives owed the spent at least the courtesy of his full attention while he spent them.

When it was done, when the two hundred had taken up their position in the narrows and the column had begun to march on without them, Jon found Daenerys watching him again from across the camp — that same careful, weighing look, held a beat too long — and understood, meeting it, that whatever had changed in her, she had just watched him do the hardest thing he had done since the war began, and he still had no idea what she was thinking as she watched it.

He almost went to her then. He almost crossed the camp and asked her plainly what had shifted between them, because the not-knowing had begun to cost him more than any answer could. But the column was marching, and two hundred men were dying behind him to make the marching possible, and there was no room left in that particular night for a question that could wait.

He let it wait. He would think, later, that he had let a great many things wait that particular winter, and that the waiting had cost him more, in the end, than the asking ever would have.